\chapter{Magnetic Resonance Imaging (MRI)}

\section*{What Is This Test?}
MRI uses powerful magnets and radio waves to create exquisitely detailed images of soft tissues. It's superior for brain, spine, joints, and organs. No radiation, but longer scan time (30-60 minutes) and not suitable for some metal implants.

\section*{Alternative Names}
MRI, magnetic resonance imaging, MR scan, nuclear magnetic resonance (NMR) imaging


\section*{Principle}
MRI uses a strong static magnetic field (1.5-7 Tesla) that aligns hydrogen protons in tissue. Radiofrequency pulses excite the protons; as they relax, they emit signals that are spatially encoded by gradient coils. T1-weighted images show anatomy; T2-weighted images show pathology. Contrast (gadolinium) alters T1 relaxation.

\section*{History}
MRI was developed independently by Paul Lauterbur and Peter Mansfield in the 1970s, earning them the Nobel Prize in 2003. The first clinical MRI scanners were installed in the early 1980s.

\section*{Why This Test Is Ordered}
Your doctor orders an MRI when they need highly detailed images of soft tissues --- the brain, spinal cord, muscles, ligaments, organs, and blood vessels. MRI is especially valuable for detecting tumors, strokes, torn ligaments, spinal disc problems, and abnormalities that other imaging cannot see well. It answers the question: what is happening in this part of your body at a level of detail that other tests cannot match?

\section*{Is This a Primary or Secondary Test}
A secondary test. Typically ordered after X-ray, CT, or ultrasound when more detailed soft tissue information is needed.

\section*{How This Test Is Performed}
You lie on a table that slides into a large tube-shaped machine. The machine uses a powerful magnet and radio waves to create detailed images of your body. The test is noisy --- you will hear loud thumping and clicking sounds, and earplugs or headphones are provided. The test takes thirty to ninety minutes depending on what is being examined. You must lie very still during the scan. You can communicate with the technician through an intercom.

\section*{What Preparation Is Needed}
Tell the MRI staff if you have:
\begin{itemize}
\item A pacemaker (unless it is specifically labeled MRI-safe)
\item Metal implants, surgical clips, or shrapnel in your body
\item A cochlear implant (artificial ear)
\item Certain aneurysm clips in your brain
\item Metal fragments in your eyes
\end{itemize}

These items can be dangerous in the strong magnetic field. You will be asked to remove all jewelry, watches, and metal objects. You will change into a hospital gown. If contrast dye (gadolinium) is to be used, your kidney function will be checked beforehand with a blood test.

\section*{Contraindications and Precautions}
MRI is contraindicated in patients with certain metal implants, pacemakers (unless MRI-conditional), cochlear implants, and some aneurysm clips. Complete the facility's MRI safety screening before the scan. Inform the team about any implants, devices, or metal in your body. Gadolinium contrast requires adequate kidney function.

\section*{Risks}
MRI uses no ionizing radiation. The strong magnetic field is safe for most patients. Risks include: discomfort from lying in the scanner for 30--90 minutes, loud noise (ear protection is provided), and (rarely) allergic reaction to gadolinium contrast dye. Patients with certain metal implants may be at risk from the magnetic field.

\section*{What Happens After the Test}
There is no recovery time from the scan itself. Routine post-scan kidney monitoring is not required for every patient; follow the agent-specific and institutional instructions. Results are typically available within one to twenty-four hours depending on the urgency.

\section*{Understanding Your Results}

\section*{Normal Results}
A normal result shows no structural abnormalities, masses, inflammation, or tissue damage in the imaged area.

\section*{Follow-up Tests}
If the MRI is abnormal, your doctor may order a contrast-enhanced MRI for better characterization, a CT scan for complementary information, a biopsy for tissue diagnosis, or a follow-up MRI to monitor changes over time.
\clearpage
